\section{Empirical Mixing}
\label{sec:empirical}

\subsection{Measurement Protocol}

To measure empirical mixing times, we run the convergence diagnostics
of G.4~\citep{b16paper} on ten parallel \recom{} chains for each
of the six study states.
We record the step count at which:
\begin{enumerate}
\item $\rhat(\PP) < 1.1$ first (and remains below 1.1 for 500 subsequent steps).
\item $\ess(\PP) > 500$ first.
\item $\ham(1) < 0.95$ (this is typically satisfied from the start).
\end{enumerate}

The maximum of these three thresholds is our empirical mixing time
$\hat{t}_{\rm mix}$.

\subsection{Results}

\begin{table}[h]
\centering
\caption{Empirical mixing times by state. $\hat{t}_{\rm mix}$ is the
         step count for all three G.4 diagnostic thresholds to be
         simultaneously satisfied. ``Bottleneck'' indicates which
         diagnostic was slowest.}
\label{tab:empirical-mixing}
\begin{tabular}{lcccc}
\toprule
State & $k$ & $n$ & $\hat{t}_{\rm mix}$ & Bottleneck \\
\midrule
NC & 14 &  2{,}195 &  2{,}000 & $\rhat$ \\
WI &  8 &  1{,}406 &  1{,}500 & $\rhat$ \\
GA & 14 &  1{,}969 &  2{,}500 & $\rhat$ \\
PA & 17 &  3{,}218 &  3{,}000 & $\rhat$ \\
TX & 38 &  5{,}265 & 10{,}000 & $\ess$ \\
CA & 52 &  8{,}057 & 20{,}000 & $\ess$ \\
\bottomrule
\end{tabular}
\end{table}

\subsection{Scaling with State Size}

The empirical mixing time scales approximately as:
\[
  \hat{t}_{\rm mix} \approx 400 \cdot k,
\]
where $k$ is the district count.
This is a much more favorable scaling than the theoretical bound
$O(n^2 \log n)$ (which scales with tract count, not district count).
The practical rate-limiter is not graph geometry but district count:
larger delegations have more districts to mix.

The relationship $\hat{t}_{\rm mix} \approx 400k$ suggests that the
relevant state space diameter, from the perspective of practical mixing,
is $\approx 400$ \recom{} steps per district rather than the
$O(n^2)$ steps suggested by spectral theory.

\subsection{Comparison to Published Practice}

Most published redistricting ensemble studies use 50{,}000--100{,}000
steps without reporting convergence diagnostics.
Under our empirical estimates, 50{,}000 steps is:
\begin{itemize}
\item $25\times$ the minimum for NC ($\hat{t}_{\rm mix} = 2{,}000$).
\item $5\times$ the minimum for TX ($\hat{t}_{\rm mix} = 10{,}000$).
\item $2.5\times$ the minimum for CA ($\hat{t}_{\rm mix} = 20{,}000$).
\end{itemize}

For most states, 50{,}000 steps is more than adequate.
The concern is not that existing studies are under-mixing, but that they
do not certify convergence.
A future study with only 5{,}000 steps for California---plausible if run
on consumer hardware---would be uncertified.

\subsection{ESS is the Binding Constraint for Large States}

For Texas and California, $\rhat < 1.1$ is achieved before $\ess > 500$.
This is because the high autocorrelation ($\hat{\rho}_1 \approx 0.92$--$0.94$)
suppresses $\ess$ even after the chains have converged to the same region
of state space.
The $\ess$ constraint is therefore the more demanding requirement for
large-delegation states.
